% Options for packages loaded elsewhere
% Options for packages loaded elsewhere
\PassOptionsToPackage{unicode}{hyperref}
\PassOptionsToPackage{hyphens}{url}
\PassOptionsToPackage{dvipsnames,svgnames,x11names}{xcolor}
%
\documentclass[
  british,
]{article}
\usepackage{xcolor}
\usepackage{amsmath,amssymb}
\setcounter{secnumdepth}{5}
\usepackage{iftex}
\ifPDFTeX
  \usepackage[T1]{fontenc}
  \usepackage[utf8]{inputenc}
  \usepackage{textcomp} % provide euro and other symbols
\else % if luatex or xetex
  \usepackage{unicode-math} % this also loads fontspec
  \defaultfontfeatures{Scale=MatchLowercase}
  \defaultfontfeatures[\rmfamily]{Ligatures=TeX,Scale=1}
\fi
\usepackage{lmodern}
\ifPDFTeX\else
  % xetex/luatex font selection
\fi
% Use upquote if available, for straight quotes in verbatim environments
\IfFileExists{upquote.sty}{\usepackage{upquote}}{}
\IfFileExists{microtype.sty}{% use microtype if available
  \usepackage[]{microtype}
  \UseMicrotypeSet[protrusion]{basicmath} % disable protrusion for tt fonts
}{}
\makeatletter
\@ifundefined{KOMAClassName}{% if non-KOMA class
  \IfFileExists{parskip.sty}{%
    \usepackage{parskip}
  }{% else
    \setlength{\parindent}{0pt}
    \setlength{\parskip}{6pt plus 2pt minus 1pt}}
}{% if KOMA class
  \KOMAoptions{parskip=half}}
\makeatother
% Make \paragraph and \subparagraph free-standing
\makeatletter
\ifx\paragraph\undefined\else
  \let\oldparagraph\paragraph
  \renewcommand{\paragraph}{
    \@ifstar
      \xxxParagraphStar
      \xxxParagraphNoStar
  }
  \newcommand{\xxxParagraphStar}[1]{\oldparagraph*{#1}\mbox{}}
  \newcommand{\xxxParagraphNoStar}[1]{\oldparagraph{#1}\mbox{}}
\fi
\ifx\subparagraph\undefined\else
  \let\oldsubparagraph\subparagraph
  \renewcommand{\subparagraph}{
    \@ifstar
      \xxxSubParagraphStar
      \xxxSubParagraphNoStar
  }
  \newcommand{\xxxSubParagraphStar}[1]{\oldsubparagraph*{#1}\mbox{}}
  \newcommand{\xxxSubParagraphNoStar}[1]{\oldsubparagraph{#1}\mbox{}}
\fi
\makeatother

\usepackage{color}
\usepackage{fancyvrb}
\newcommand{\VerbBar}{|}
\newcommand{\VERB}{\Verb[commandchars=\\\{\}]}
\DefineVerbatimEnvironment{Highlighting}{Verbatim}{commandchars=\\\{\}}
% Add ',fontsize=\small' for more characters per line
\usepackage{framed}
\definecolor{shadecolor}{RGB}{241,243,245}
\newenvironment{Shaded}{\begin{snugshade}}{\end{snugshade}}
\newcommand{\AlertTok}[1]{\textcolor[rgb]{0.68,0.00,0.00}{#1}}
\newcommand{\AnnotationTok}[1]{\textcolor[rgb]{0.37,0.37,0.37}{#1}}
\newcommand{\AttributeTok}[1]{\textcolor[rgb]{0.40,0.45,0.13}{#1}}
\newcommand{\BaseNTok}[1]{\textcolor[rgb]{0.68,0.00,0.00}{#1}}
\newcommand{\BuiltInTok}[1]{\textcolor[rgb]{0.00,0.23,0.31}{#1}}
\newcommand{\CharTok}[1]{\textcolor[rgb]{0.13,0.47,0.30}{#1}}
\newcommand{\CommentTok}[1]{\textcolor[rgb]{0.37,0.37,0.37}{#1}}
\newcommand{\CommentVarTok}[1]{\textcolor[rgb]{0.37,0.37,0.37}{\textit{#1}}}
\newcommand{\ConstantTok}[1]{\textcolor[rgb]{0.56,0.35,0.01}{#1}}
\newcommand{\ControlFlowTok}[1]{\textcolor[rgb]{0.00,0.23,0.31}{\textbf{#1}}}
\newcommand{\DataTypeTok}[1]{\textcolor[rgb]{0.68,0.00,0.00}{#1}}
\newcommand{\DecValTok}[1]{\textcolor[rgb]{0.68,0.00,0.00}{#1}}
\newcommand{\DocumentationTok}[1]{\textcolor[rgb]{0.37,0.37,0.37}{\textit{#1}}}
\newcommand{\ErrorTok}[1]{\textcolor[rgb]{0.68,0.00,0.00}{#1}}
\newcommand{\ExtensionTok}[1]{\textcolor[rgb]{0.00,0.23,0.31}{#1}}
\newcommand{\FloatTok}[1]{\textcolor[rgb]{0.68,0.00,0.00}{#1}}
\newcommand{\FunctionTok}[1]{\textcolor[rgb]{0.28,0.35,0.67}{#1}}
\newcommand{\ImportTok}[1]{\textcolor[rgb]{0.00,0.46,0.62}{#1}}
\newcommand{\InformationTok}[1]{\textcolor[rgb]{0.37,0.37,0.37}{#1}}
\newcommand{\KeywordTok}[1]{\textcolor[rgb]{0.00,0.23,0.31}{\textbf{#1}}}
\newcommand{\NormalTok}[1]{\textcolor[rgb]{0.00,0.23,0.31}{#1}}
\newcommand{\OperatorTok}[1]{\textcolor[rgb]{0.37,0.37,0.37}{#1}}
\newcommand{\OtherTok}[1]{\textcolor[rgb]{0.00,0.23,0.31}{#1}}
\newcommand{\PreprocessorTok}[1]{\textcolor[rgb]{0.68,0.00,0.00}{#1}}
\newcommand{\RegionMarkerTok}[1]{\textcolor[rgb]{0.00,0.23,0.31}{#1}}
\newcommand{\SpecialCharTok}[1]{\textcolor[rgb]{0.37,0.37,0.37}{#1}}
\newcommand{\SpecialStringTok}[1]{\textcolor[rgb]{0.13,0.47,0.30}{#1}}
\newcommand{\StringTok}[1]{\textcolor[rgb]{0.13,0.47,0.30}{#1}}
\newcommand{\VariableTok}[1]{\textcolor[rgb]{0.07,0.07,0.07}{#1}}
\newcommand{\VerbatimStringTok}[1]{\textcolor[rgb]{0.13,0.47,0.30}{#1}}
\newcommand{\WarningTok}[1]{\textcolor[rgb]{0.37,0.37,0.37}{\textit{#1}}}

\usepackage{longtable,booktabs,array}
\newcounter{none} % for unnumbered tables
\usepackage{calc} % for calculating minipage widths
% Correct order of tables after \paragraph or \subparagraph
\usepackage{etoolbox}
\makeatletter
\patchcmd\longtable{\par}{\if@noskipsec\mbox{}\fi\par}{}{}
\makeatother
% Allow footnotes in longtable head/foot
\IfFileExists{footnotehyper.sty}{\usepackage{footnotehyper}}{\usepackage{footnote}}
\makesavenoteenv{longtable}
\usepackage{graphicx}
\makeatletter
\newsavebox\pandoc@box
\newcommand*\pandocbounded[1]{% scales image to fit in text height/width
  \sbox\pandoc@box{#1}%
  \Gscale@div\@tempa{\textheight}{\dimexpr\ht\pandoc@box+\dp\pandoc@box\relax}%
  \Gscale@div\@tempb{\linewidth}{\wd\pandoc@box}%
  \ifdim\@tempb\p@<\@tempa\p@\let\@tempa\@tempb\fi% select the smaller of both
  \ifdim\@tempa\p@<\p@\scalebox{\@tempa}{\usebox\pandoc@box}%
  \else\usebox{\pandoc@box}%
  \fi%
}
% Set default figure placement to htbp
\def\fps@figure{htbp}
\makeatother


% definitions for citeproc citations
\NewDocumentCommand\citeproctext{}{}
\NewDocumentCommand\citeproc{mm}{%
  \begingroup\def\citeproctext{#2}\cite{#1}\endgroup}
\makeatletter
 % allow citations to break across lines
 \let\@cite@ofmt\@firstofone
 % avoid brackets around text for \cite:
 \def\@biblabel#1{}
 \def\@cite#1#2{{#1\if@tempswa , #2\fi}}
\makeatother
\newlength{\cslhangindent}
\setlength{\cslhangindent}{1.5em}
\newlength{\csllabelwidth}
\setlength{\csllabelwidth}{3em}
\newenvironment{CSLReferences}[2] % #1 hanging-indent, #2 entry-spacing
 {\begin{list}{}{%
  \setlength{\itemindent}{0pt}
  \setlength{\leftmargin}{0pt}
  \setlength{\parsep}{0pt}
  % turn on hanging indent if param 1 is 1
  \ifodd #1
   \setlength{\leftmargin}{\cslhangindent}
   \setlength{\itemindent}{-1\cslhangindent}
  \fi
  % set entry spacing
  \setlength{\itemsep}{#2\baselineskip}}}
 {\end{list}}
\usepackage{calc}
\newcommand{\CSLBlock}[1]{\hfill\break\parbox[t]{\linewidth}{\strut\ignorespaces#1\strut}}
\newcommand{\CSLLeftMargin}[1]{\parbox[t]{\csllabelwidth}{\strut#1\strut}}
\newcommand{\CSLRightInline}[1]{\parbox[t]{\linewidth - \csllabelwidth}{\strut#1\strut}}
\newcommand{\CSLIndent}[1]{\hspace{\cslhangindent}#1}

\ifLuaTeX
\usepackage[bidi=basic,shorthands=off]{babel}
\else
\usepackage[bidi=default,shorthands=off]{babel}
\fi
\ifLuaTeX
  \usepackage{selnolig} % disable illegal ligatures
\fi


\setlength{\emergencystretch}{3em} % prevent overfull lines

\providecommand{\tightlist}{%
  \setlength{\itemsep}{0pt}\setlength{\parskip}{0pt}}



 


% Get Math Fonts
\usepackage{fontspec}
\usepackage{unicode-math}
\setmathfont{Latin Modern Math}
\setmathfont[range={\mathscr,\mathbfscr,\mathbb}]{XITSMath-Regular}
[    Extension = .otf,
      BoldFont = XITSMath-Bold
]
\makeatletter
\@ifpackageloaded{caption}{}{\usepackage{caption}}
\AtBeginDocument{%
\ifdefined\contentsname
  \renewcommand*\contentsname{Table of contents}
\else
  \newcommand\contentsname{Table of contents}
\fi
\ifdefined\listfigurename
  \renewcommand*\listfigurename{List of Figures}
\else
  \newcommand\listfigurename{List of Figures}
\fi
\ifdefined\listtablename
  \renewcommand*\listtablename{List of Tables}
\else
  \newcommand\listtablename{List of Tables}
\fi
\ifdefined\figurename
  \renewcommand*\figurename{Figure}
\else
  \newcommand\figurename{Figure}
\fi
\ifdefined\tablename
  \renewcommand*\tablename{Table}
\else
  \newcommand\tablename{Table}
\fi
}
\@ifpackageloaded{float}{}{\usepackage{float}}
\floatstyle{ruled}
\@ifundefined{c@chapter}{\newfloat{codelisting}{h}{lop}}{\newfloat{codelisting}{h}{lop}[chapter]}
\floatname{codelisting}{Listing}
\newcommand*\listoflistings{\listof{codelisting}{List of Listings}}
\makeatother
\makeatletter
\makeatother
\makeatletter
\@ifpackageloaded{caption}{}{\usepackage{caption}}
\@ifpackageloaded{subcaption}{}{\usepackage{subcaption}}
\makeatother
\usepackage{bookmark}
\IfFileExists{xurl.sty}{\usepackage{xurl}}{} % add URL line breaks if available
\urlstyle{same}
% fallback for those not using the hyperref driver hyperxmp:
\makeatletter
\@ifundefined{xmpquote}{\newcommand{\xmpquote}[1]{#1}}{}
\makeatother
\hypersetup{
  pdftitle={Designing an Alternative Approach to Mathematical Research},
  pdfauthor={Kai Prince},
  pdflang={en-GB},
  pdfkeywords={\xmpquote{Mathematical Research}, \xmpquote{Open
Research}, \xmpquote{Open Education}, \xmpquote{Universal Design for
Learning}, \xmpquote{Disability Inclusion}},
  colorlinks=true,
  linkcolor={blue},
  filecolor={Maroon},
  citecolor={Blue},
  urlcolor={Blue},
  pdfcreator={LaTeX via pandoc}}


\title{Designing an Alternative Approach to Mathematical Research}
\author{Kai Prince}
\date{2026-09-27}
\begin{document}
\maketitle


{[}Book Project{]}

{[}Check Journal-format Thesis requirements{]}

\section{Introduction}\label{introduction}

Concept: Designing a Quarto Extension to generate an Open Education
Website to maximise the accessibility of Open Research projects.

{[}Insert OpenMath motivations{]}

{[}Insert Disability Inclusion in Research motivations{]}

\section{Open Research}\label{open-research}

The ``Position Statement on Open Research'' at The University of
Manchester (\citeproc{ref-uom2025OpenResPosition}{2025c}) covers what
Open Research is, why it is important and its principles:

\begin{itemize}
\tightlist
\item
  Pre-registration of research;
\item
  Transparency in research methodology;
\item
  Public availability and reusability of research data and analysis
  code;
\item
  Public accessibility and transparency of research communication;
\item
  Data described according to FAIR Data Principles, ensuring that it is
  Findable, Accessible Interoperable and Reusable and publicly available
  where possible;
\item
  Using web-based tools to facilitate collaboration.
\end{itemize}

\subsection{UoM Policies}\label{uom-policies}

The University of Manchester (\citeproc{ref-uom2025PubPolicy}{2025a})
details state that researchers must be open and transparent when
conducting and communicating their research in terms of:

\begin{enumerate}
\def\labelenumi{\arabic{enumi}.}
\tightlist
\item
  the disclosure of any conflicts of interest;
\item
  the reporting of research data collection methods;
\item
  the analysis and interpretation of data;
\item
  making all research findings widely available (including sharing
  negative results as appropriate) by publishing as open access under a
  Creative Commons Attribution (CC BY) license;
\item
  disseminating research in a way that will maximise engagement and
  secondary use, with students obtaining prior permission from their
  supervisor; and
\item
  promoting public engagement/involvement in research.
\end{enumerate}

The University of Manchester
(\citeproc{ref-uom2024ResDataMngmtPolicy}{2024}) that research data must
be stored in an ``accurate, complete unadulterated and reliable manner''
as well as meeting FAIR data principles, where the data must be made
interoperable whenever possible, and ``made available in a timely way
with as few restrictions as possible''.

The University of Manchester (\citeproc{ref-uom2025RespMetrics}{2025b})
endorse five principles of responsible metrics:

\begin{itemize}
\tightlist
\item
  Robustness
\item
  Humility
\item
  Transparency
\item
  Diversity
\item
  Reflexivity
\end{itemize}

\section{Mathematical Research}\label{mathematical-research}

\subsection{Researcher Lifecycle}\label{researcher-lifecycle}

{[}Conceptual overlaps between Open Research \& Open Education{]}

\subsection{Mathematical Research
Methodologies}\label{mathematical-research-methodologies}

\subsection{Knowledge Exchange
Initiatives}\label{knowledge-exchange-initiatives}

\subsection{Research Outputs}\label{research-outputs}

{[}Papers{]}

{[}Presentations{]}

{[}Posters{]}

{[}Social Media Content{]}

{[}Public Outreach{]}

\subsection{Proof Verification}\label{proof-verification}

{[}Rocq/Lean{]}

\subsection{Open Research in
Mathematics}\label{open-research-in-mathematics}

{[}Accessibility{]}

{[}Findability{]}

{[}Interoperability{]}

{[}Reusability{]}

{[}Below is from Knowledge Graphs YouTube Lecture 1.7{]}

Linked Data Principles:

\begin{enumerate}
\def\labelenumi{\arabic{enumi}.}
\tightlist
\item
  Use URIs as names for things.
\item
  Use HTTP URIs so that people can look up those names.
\item
  When someone looks up a URI, provide useful information, using the
  standards (Resource Description Framework (RDF), SPARQL).
\item
  Include links to other URIs, so they can discover more things.
\end{enumerate}

Tim Berners-Lee's 5-Star Criteria for Linked Open Data:

\begin{enumerate}
\def\labelenumi{\arabic{enumi}.}
\tightlist
\item
  Available on the Web (whatever format) but with an \textbf{open
  license}, to be Open Data.
\item
  Available as \textbf{machine-readable structured data} (e.g., excel
  instead of image scan of a table).
\item
  As (2) plus \textbf{non-proprietary format} (e.g., CSV instead of
  Excel).
\item
  All of the above plus: use \textbf{open standards from W3C} (RDF and
  SPARQL) to identify things so that people can point at your stuff.
\item
  All the above plus: \textbf{link your data to other people's data} to
  provide context.
\end{enumerate}

\subsection{Ethical AI Research
Practices}\label{ethical-ai-research-practices}

\section{Authoring}\label{authoring}

{[}Format: What is it; What aspects they already cover; Summarise
remaining actions{]}

\subsection{Executable Papers \& Open Education
Resources}\label{executable-papers-open-education-resources}

{[}Open Access \& Open Education{]}

{[}Universal Design for Learning{]}

{[}Executable Papers{]}

\subsection{Platforms}\label{platforms}

\subsubsection{Overleaf}\label{overleaf}

\subsubsection{Obsidian}\label{obsidian}

\subsubsection{Quarto \& Quarto
Extensions}\label{quarto-quarto-extensions}

{[}Journal Article Templates{]}

\section{Backend Development}\label{backend-development}

{[}Approach - Algorithmic wherever possible{]}

\subsection{Rendering Workflow \& Knowledge
Graphs}\label{rendering-workflow-knowledge-graphs}

Data modelling process:

\begin{enumerate}
\def\labelenumi{\arabic{enumi}.}
\tightlist
\item
  Understand the domain and define specific use cases (questions) for
  the application
\item
  Develop the initial graph data model:

  \begin{enumerate}
  \def\labelenumii{\arabic{enumii}.}
  \tightlist
  \item
    Model the nodes (entities).

    \begin{itemize}
    \tightlist
    \item
      Choose sematically orthogonal labels and avoid class hierarchies
      in labels or using more than 4 labels.
    \end{itemize}
  \item
    Model the relationships between nodes.
  \end{enumerate}
\item
  Test the use cases against the initial data model.
\item
  Create the graph (instance model) with test data using Cypher.
\item
  Test the use cases, including performance against the graph.
\item
  Refactor (improve) the graph data model due to a change in the key use
  cases or for performance reasons.

  \begin{enumerate}
  \def\labelenumii{\arabic{enumii}.}
  \tightlist
  \item
    Design the new data model.
  \item
    Write Cypher code to transform the existing graph to implement the
    new data model.
  \item
    Retest all use cases, possibly with updated Cypher code.
  \end{enumerate}
\item
  Implement the refactoring on the graph and retest using Cypher.

  \begin{enumerate}
  \def\labelenumii{\arabic{enumii}.}
  \tightlist
  \item
    Identify which use cases are affected by the refactor.
  \item
    Rewrite any queries that can take advantage of the refactoring.
  \item
    Test all queries affected by the refactor to ensure they return the
    same results as before the refactoring.
  \item
    PROFILE queries to determine if refactoring is a benefit.
  \end{enumerate}
\end{enumerate}

Use cases:

\begin{enumerate}
\def\labelenumi{\arabic{enumi}.}
\tightlist
\item
  What are the \textbf{sources} \texttt{cited} for each \textbf{result}?

  \begin{enumerate}
  \def\labelenumii{\arabic{enumii}.}
  \tightlist
  \item
    Retrieve a result by its \textbf{embedding vector}.
  \item
    Trace the CITED relationships to the sources.
  \item
    Return the \textbf{citations} of the sources.
  \end{enumerate}
\item
  What is the \textbf{result} that \texttt{defines} a specific
  \textbf{object}?

  \begin{enumerate}
  \def\labelenumii{\arabic{enumii}.}
  \tightlist
  \item
    Retrieve the object by its \textbf{name}.
  \item
    Trace the DEFINE relationship to the result.
  \item
    Return the \textbf{description} and \textbf{href} of the result.
  \end{enumerate}
\item
  What \textbf{objects} are \texttt{required} for a \textbf{result}?

  \begin{enumerate}
  \def\labelenumii{\arabic{enumii}.}
  \tightlist
  \item
    Retrieve the result by its \textbf{embedding vector}.
  \item
    Trace the REQUIRE relationship to the objects.
  \item
    Return the \textbf{names} and \textbf{notation} of the objects.
  \end{enumerate}
\item
  What objects \texttt{compose} a specific object?

  \begin{enumerate}
  \def\labelenumii{\arabic{enumii}.}
  \tightlist
  \item
    Retrieve the specific object by its \textbf{name}.
  \item
    Trace the composed DEFINE and REQUIRE relationships to objects.
  \item
    Return the \textbf{names} of the objects.
  \end{enumerate}
\item
  What \textbf{properties} \texttt{apply} to a specific \textbf{object}?

  \begin{enumerate}
  \def\labelenumii{\arabic{enumii}.}
  \tightlist
  \item
    Retrieve the specific object by its \textbf{name}.
  \item
    Trace the composed DEFINE, REQUIRE and APPLY relationships to
    objects.
  \item
    Return the \textbf{names} of the objects.
  \end{enumerate}
\item
  What nodes \texttt{originated} from a specific \textbf{document}?

  \begin{enumerate}
  \def\labelenumii{\arabic{enumii}.}
  \tightlist
  \item
    Retrieve the document by its \textbf{path}.
  \item
    Trace the ORIGIN relationships to the nodes.
  \item
    Return the key values and labels of the nodes.
  \end{enumerate}
\end{enumerate}

Nodes:

\begin{itemize}
\tightlist
\item
  SOURCE

  \begin{itemize}
  \tightlist
  \item
    {[}{[}referenceID{]}{]}
  \item
    citation
  \end{itemize}
\item
  DOCUMENT

  \begin{itemize}
  \tightlist
  \item
    {[}{[}path{]}{]}
  \end{itemize}
\item
  OBJECT:

  \begin{itemize}
  \tightlist
  \item
    {[}{[}identifier{]}{]}
  \item
    notation
  \end{itemize}
\item
  RESULT

  \begin{itemize}
  \tightlist
  \item
    {[}{[}embeddingVec{]}{]}
  \item
    {[}abstractionVec{]}
  \item
    desc
  \item
    href
  \end{itemize}
\end{itemize}

Intermediate Nodes:

\begin{itemize}
\tightlist
\item
  RESULT

  \begin{itemize}
  \tightlist
  \item
    create applied result with abstraction relationship to the original
    result unless the original result is sufficiently abstract
  \end{itemize}
\end{itemize}

Relationships:

\begin{itemize}
\tightlist
\item
  CITE
\item
  PROPERTY
\item
  ABSTRACTION
\item
  APPLIED
\end{itemize}

{[}SQ4R{]}

Quarto processing:

\begin{enumerate}
\def\labelenumi{\arabic{enumi}.}
\setcounter{enumi}{-1}
\tightlist
\item
  Source Documents
\item
  Pre-render
\item
  Render Filters

  \begin{itemize}
  \tightlist
  \item
    Pandoc Structure:

    \begin{itemize}
    \tightlist
    \item
      Metadata
    \item
      Blocks

      \begin{itemize}
      \tightlist
      \item
        Blocks
      \item
        Inlines
      \end{itemize}
    \end{itemize}
  \end{itemize}
\item
  Post-render
\item
  Output Documents
\end{enumerate}

Embedding RAG:

\begin{enumerate}
\def\labelenumi{\arabic{enumi}.}
\setcounter{enumi}{-1}
\tightlist
\item
  Source Documents
\item
  Data Ingestion
\item
  (Unstructured) Text Chunks
\item
  Embeddings (using an Embedding Model)
\item
  Vector Database

  \begin{itemize}
  \tightlist
  \item
    Semantic-aware
  \end{itemize}
\end{enumerate}

GraphRAG:

\begin{enumerate}
\def\labelenumi{\arabic{enumi}.}
\setcounter{enumi}{-1}
\tightlist
\item
  Source Documents
\item
  Chunks
\item
  Knowledge Graph Generation

  \begin{itemize}
  \tightlist
  \item
    Entity Extraction (using an LLM)
  \item
    Relationship Extraction (using an LLM)
  \end{itemize}
\item
  Community Detection (using an LLM)
\item
  Hierarchical Community Structures (using an LLM)
\item
  Community Levels - Root, Low, or High (using an LLM)
\item
  Generate Community Summaries (using an LLM)
\item
  Vector Database (using an Embedding Model)

  \begin{itemize}
  \tightlist
  \item
    Context-aware
  \end{itemize}
\end{enumerate}

Real-Time Knowledge Graphs (Graphiti for dynamic data management):

\begin{enumerate}
\def\labelenumi{\arabic{enumi}.}
\setcounter{enumi}{-1}
\tightlist
\item
  Source Documents
\item
  Chunks
\item
  Knowledge Graphs Construction
\end{enumerate}

{\def\LTcaptype{none} % do not increment counter
\begin{longtable}[]{@{}
  >{\raggedright\arraybackslash}p{(\linewidth - 8\tabcolsep) * \real{0.2000}}
  >{\raggedright\arraybackslash}p{(\linewidth - 8\tabcolsep) * \real{0.2000}}
  >{\raggedright\arraybackslash}p{(\linewidth - 8\tabcolsep) * \real{0.2000}}
  >{\raggedright\arraybackslash}p{(\linewidth - 8\tabcolsep) * \real{0.2000}}
  >{\raggedright\arraybackslash}p{(\linewidth - 8\tabcolsep) * \real{0.2000}}@{}}
\toprule\noalign{}
\begin{minipage}[b]{\linewidth}\raggedright
Stage
\end{minipage} & \begin{minipage}[b]{\linewidth}\raggedright
Quarto
\end{minipage} & \begin{minipage}[b]{\linewidth}\raggedright
Embedding RAG
\end{minipage} & \begin{minipage}[b]{\linewidth}\raggedright
GraphRAG
\end{minipage} & \begin{minipage}[b]{\linewidth}\raggedright
Temporal RAG
\end{minipage} \\
\midrule\noalign{}
\endhead
\bottomrule\noalign{}
\endlastfoot
Input & Markdown Documents & & & \\
Output & Formatted Documents & Semantic-aware Vector Database &
Context-aware Vector Database & Temporal-aware Vector Database \\
\end{longtable}
}

{[}Directory specific{]}

\subsection{Data Processing}\label{data-processing}

\subsubsection{Data Cleaning}\label{data-cleaning}

{[}Automatic Abbreviation Detection{]}

{[}Override/Example processing{]}

\subsubsection{Mathematical Data}\label{mathematical-data}

\subsubsection{Mathematical Proofs}\label{mathematical-proofs}

\subsubsection{Terminology, Definitions \&
Conclusions}\label{terminology-definitions-conclusions}

{[}Definition/Div/Paragraph processing, e.g., emph text = defined
term{]}

Basic Definition Schema:

\begin{itemize}
\tightlist
\item
  name: emph text
\item
  container\_notation
\item
  element\_notation
\end{itemize}

Definition Relationships:

\begin{itemize}
\tightlist
\item
  REQUIRES -\textgreater{} Definition
\item
  HAS\_PROPERTY -\textgreater{} Definition
\item
  SUBCATEGORY\_OF -\textgreater{} Definition
\item
  CONTAINS -\textgreater{} Definition
\end{itemize}

{[}Synonym Sets{]}

{[}Language-Independent URIs for Designator and Designatum, e.g.,
Wikidata Identifier{]}

Want to get information about ``the apple'' from Wikidata:

\begin{Shaded}
\begin{Highlighting}[]
\NormalTok{HTTP GET request}
\NormalTok{Accept Header: text/html}
\NormalTok{https://wikidata.org/entity/Q89}
\end{Highlighting}
\end{Shaded}

URI represents Designatum

\begin{Shaded}
\begin{Highlighting}[]
\NormalTok{HTTP/2 303 See Other}
\NormalTok{https://wikidata.org/wiki/Q89}
\end{Highlighting}
\end{Shaded}

URI represents Designator

\begin{Shaded}
\begin{Highlighting}[]
\NormalTok{HTTP GET request}
\NormalTok{Accept Header: text/html}
\NormalTok{https://wikidata.org/wiki/Q89}
\end{Highlighting}
\end{Shaded}

Gets HTML Document

Want to get machine understandable information about ``the apple'' from
Wikidata:

\begin{Shaded}
\begin{Highlighting}[]
\NormalTok{HTTP GET request}
\NormalTok{Accept Header: text/turtle}
\NormalTok{https://wikidata.org/entity/Q89}
\end{Highlighting}
\end{Shaded}

URI represents Designatum

\begin{Shaded}
\begin{Highlighting}[]
\NormalTok{HTTP/2 303 See Other}
\NormalTok{https://wikidata.org/wiki/Special:EntityData/Q89.ttl}
\end{Highlighting}
\end{Shaded}

URI represents Designator

\begin{Shaded}
\begin{Highlighting}[]
\NormalTok{HTTP GET request}
\NormalTok{Accept Header: text/turtle}
\NormalTok{https://wikidata.org/wiki/Special:EntityData/Q89.ttl}
\end{Highlighting}
\end{Shaded}

Gets HTML Document RDF Document (Turtle is a way to encode RDF)

\paragraph{Anticipating International
Translations}\label{anticipating-international-translations}

\paragraph{Addressing Biases}\label{addressing-biases}

{[}Decolonising the Curriculum{]}

\subsection{Data Management}\label{data-management}

\subsubsection{Data Processing}\label{data-processing-1}

\paragraph{Dependency Identification \&
Sorting}\label{dependency-identification-sorting}

\paragraph{Internal Reference
Management}\label{internal-reference-management}

\paragraph{External Reference
Management}\label{external-reference-management}

\paragraph{Cross Referencing \&
Embedding}\label{cross-referencing-embedding}

\paragraph{Version Tracking}\label{version-tracking}

\subsubsection{Error Handling}\label{error-handling}

\subsection{Document Rendering, Progressive JSON \&
CSS}\label{document-rendering-progressive-json-css}

\paragraph{LaTeX, MathJax \& Presentation
MathML}\label{latex-mathjax-presentation-mathml}

{[}Math Environment HTML Promises{]}

\paragraph{Alt-Text, MathSpeak \& Context
MathML}\label{alt-text-mathspeak-context-mathml}

\section{Frontend Development}\label{frontend-development}

\subsection{PDF Output}\label{pdf-output}

\subsection{Web Output}\label{web-output}

\subsubsection{Progressive Web App}\label{progressive-web-app}

\subsubsection{Accessibility Options and
Features}\label{accessibility-options-and-features}

\subsubsection{Linked Knowledge}\label{linked-knowledge}

\subsubsection{Semantic Search}\label{semantic-search}

{[}Jina.AI?{]}

\subsubsection{Customisable Notation \&
Terminology}\label{customisable-notation-terminology}

\subsubsection{Interactive Proofs}\label{interactive-proofs}

\subsubsection{Personalised \& Varied
Learning}\label{personalised-varied-learning}

\paragraph{Goal-Setting \& Progress
Monitoring}\label{goal-setting-progress-monitoring}

\paragraph{Assisted Note-Making}\label{assisted-note-making}

\paragraph{Playful Learning}\label{playful-learning}

\section{Discussion}\label{discussion}

\subsection{Copyright Considerations \& the Importance of
Individuality}\label{copyright-considerations-the-importance-of-individuality}

\subsection{Research Implications}\label{research-implications}

{[}Research Progression through Blog writing{]}

\subsection{Teaching Implications}\label{teaching-implications}

\subsection{Interdisciplinary
Contexts}\label{interdisciplinary-contexts}

\protect\phantomsection\label{refs}
\begin{CSLReferences}{1}{0}
\bibitem[\citeproctext]{ref-uom2024ResDataMngmtPolicy}
The University of Manchester (2024). \emph{Research data management
policy}. Available at:
\url{https://documents.manchester.ac.uk/DocuInfo.aspx?DocID=33802}

\bibitem[\citeproctext]{ref-uom2025PubPolicy}
The University of Manchester (2025a). \emph{Publications policy}.
Available at:
\url{https://documents.manchester.ac.uk/DocuInfo.aspx?DocID=76067}

\bibitem[\citeproctext]{ref-uom2025RespMetrics}
The University of Manchester (2025b). \emph{Responsible metrics
statement}. Available at:
\url{https://documents.manchester.ac.uk/DocuInfo.aspx?DocID=75142}

\bibitem[\citeproctext]{ref-uom2025OpenResPosition}
The University of Manchester (2025c). \emph{University of manchester
position statement on open research}. Available at:
\url{https://documents.manchester.ac.uk/DocuInfo.aspx?DocID=55136}

\end{CSLReferences}




\end{document}
